%%BLOCK after="A nonzero coefficient theory that forgets components"
\begin{dcomment}{development}{p5:Mais pour les dérivateurs utiles en pratique}{Five coefficient theories and their classes of equivalences}
The localizer $W_{\mathbb D}$ of this chapter uses constant coefficients only. Chapter~I has finer classes, defined with locally constant coefficients: the functors $f$ for which $f^*$ is fully faithful on locally constant objects (1-complete), and those for which it is an equivalence there (2-complete); write $W^1_{\mathbb D}$ and $W^{\mathrm{lc}}_{\mathbb D}$. The second is again a fundamental localizer, for every derivator (commentary on Chapter~I, ms.~100); the first is not. For the usual coefficient theories the three classes are as follows; $W_r$ denotes the functors inducing a bijection on $\pi_0$ and isomorphisms on $\pi_i$ for $1\leq i\leq r$ at all base points, and ``fibres'' are homotopy fibres of the map of classifying spaces.

\begin{center}
\small
\begin{tabular}{p{.19\textwidth}p{.2\textwidth}p{.24\textwidth}p{.27\textwidth}}
\hline
Coefficients & $W_{\mathbb D}$ & $W^1_{\mathbb D}$ & $W^{\mathrm{lc}}_{\mathbb D}$\\
\hline
sets & $W_0$ & $\pi_0$ bijective, $\pi_1$ onto & $W_1$\\
$r$-truncated types & $W_r$ & fibres $r$-connected & $W_{r+1}$\\
spaces & $W_\infty$ & $W_\infty$ & $W_\infty$\\
complexes of $k$-modules, $k\neq0$ & $k$-homology equivalences & fibres $k$-acyclic & $\pi_0$, $\pi_1$ isomorphic and universal covers $k$-homology equivalent\\
the same for $k=\mathbf Q$ & rational homology equivalences & fibres rationally acyclic & $\pi_0$, $\pi_1$ isomorphic, $\pi_j\otimes\mathbf Q$ isomorphic for $j\geq2$\\
spectra & integral homology equivalences & fibres integrally acyclic & $W_\infty$\\
\hline
\end{tabular}
\end{center}

In the third column the classes of functors preserving cohomology, respectively homology, with locally constant coefficients coincide with $W^1_{\mathbb D}$ for these theories. The derivations of the table are not reproduced here. Three consequences bear on the questions of these pages.

(1) \emph{Component detection.} The inclusion $W^{\mathrm{lc}}_{\mathbb D}\subset W_0$ holds if and only if $\mathbb D(e)$ is not equivalent to the terminal category, whereas $W_{\mathbb D}\subset W_0$ holds if and only if $\mathbb D(e)$ is not equivalent to a preorder; in a nontrivial preorder such as the lattice of the preceding example, $W_{\mathbb D}$ consists of the functors $A\to B$ with $A$ and $B$ both empty or both nonempty. So condition L~a) for $W_{\mathbb D}$ needs two distinct parallel arrows in $\mathbb D(e)$, and for $W^{\mathrm{lc}}_{\mathbb D}$ it needs only nontriviality.

(2) \emph{The axioms Der 3 c of Chapter~I}, which ask that $W_{\mathbb D}$ be contained in the finer classes, hold for spaces and fail in every other line of the table.

(3) \emph{Which rings see the whole homotopy type.} The last column for complexes of $k$-modules is $W_\infty$ exactly when no prime number is invertible in $k$; in particular for $k=\mathbf Z$. This follows from the next proposition.

\textbf{Proposition (arithmetic of acyclicity).}\ours
\textit{Let $C$ be a complex of free abelian groups and $k$ a nonzero ring. Then $C\otimes k$ is acyclic if and only if $C\otimes\mathbf F_p$ is acyclic for every prime $p$ not invertible in $k$ and, when $k$ has characteristic zero, $C\otimes\mathbf Q$ is acyclic.}

\textit{Proof.}
By the universal coefficient sequence, $C\otimes A$ is acyclic if and only if $H\otimes A=0$ and $\operatorname{Tor}(H,A)=0$ for every homology group $H$ of $C$. For $A=\mathbf F_p$ this says that $p$ is bijective on $H$, and for $A=\mathbf Q$ that $H$ is a torsion group.

Suppose $H\otimes k=0=\operatorname{Tor}(H,k)$, and let $p$ be not invertible in $k$. Then $k/pk$ is a nonzero $\mathbf F_p$-vector space, so $\mathbf F_p$ is a direct summand of the abelian group $k/pk$. The group $H\otimes pk$ is a quotient of $H\otimes k$, hence zero, and the exact sequence $\operatorname{Tor}(H,k)\to\operatorname{Tor}(H,k/pk)\to H\otimes pk$ gives $\operatorname{Tor}(H,k/pk)=0$; also $H\otimes k/pk=0$. So $H\otimes\mathbf F_p=0=\operatorname{Tor}(H,\mathbf F_p)$. If $k$ has characteristic zero, let $T$ be its torsion subgroup; $k/T$ is torsion-free and contains $\mathbf Z$. For $H'=H/H_{\mathrm{tors}}$ the group $H'\otimes k/T$ is a quotient of $H\otimes k$, hence zero, and it contains $H'\otimes\mathbf Z$ because $H'$ is flat. So $H'=0$.

Conversely suppose $p$ is bijective on $H$ for every $p$ not invertible in $k$. If $k$ has characteristic $n>0$, the primes not invertible in $k$ are those dividing $n$; then $n$ is bijective on $H$, the derived tensor product of $H$ with $\mathbf Z/n$ vanishes, and so does that with the $\mathbf Z/n$-module $k$. If $k$ has characteristic zero and $H$ is a torsion group, every element of $H$ has an order $m$ which is a product of primes invertible in $k$; then $H\otimes k=0$, and $\operatorname{Tor}(H,k)$, a filtered colimit of groups of $m$-torsion elements of $k$ for such $m$, vanishes.\hfill$\square$

\textit{Proof of (3) from the table.}
If no prime is invertible in $k$, then $k$ has characteristic zero, and by the proposition a $k$-homology equivalence is one for all $\mathbf F_p$ and for $\mathbf Q$, hence an integral one: the homology of its cone is torsion with every prime acting bijectively. A map inducing isomorphisms on $\pi_0$ and $\pi_1$ and integral homology isomorphisms of universal covers is a weak equivalence. If a prime $q$ is invertible in $k$, the space $K(\mathbf Z/q,2)$ is simply connected and its reduced integral homology is $q$-primary torsion, so by the proposition it is $k$-acyclic; the map $K(\mathbf Z/q,2)\to e$ is then in the class of the last column and not in $W_\infty$.\hfill$\square$

A localizer $W$ is realized as $W_{\mathbb D}$ for a derivator of localized spaces whenever $W$ is accessible (commentary at ms.~3); this applies to $W^{\mathrm{lc}}_{\mathbb D}$ when it is accessible, and then the passage from $\mathbb D$ to that derivator replaces constant by locally constant coefficients. For sets it produces the tower $W_0\supset W_1\supset W_2\supset\cdots$ with intersection $W_\infty$, each stage being the theory of locally constant coefficients of the previous one.
\end{dcomment}
